\documentclass[12pt]{article}
\usepackage[utf8]{inputenc}
\usepackage[T1]{fontenc}
\usepackage{hyperref}
\title{Texistepec Popoluca Language Resources for K-12 Curriculum Development}
\author{Compiled by Claude Sonnet 5.5 (claude-sonnet-5-5) for Andr\'as Kornai}
\date{}
\begin{document}
\maketitle
\section{Introduction}
Texistepec Popoluca (ISO 639-3 poq; native name W\"a\"a 'oot) is a Mixe-Zoquean language of Veracruz, Mexico, critically endangered. Speaker counts in sources vary widely, from 1 to 370. Documentation is mostly the work of S{\o}ren Wichmann. Do not confuse it with Sierra (Highland) or Sayula Popoluca.
\section{Primary Educational Resources}
No primers or school materials were found in search results.
\section{Linguistic and Reference Materials}
\begin{itemize}
\item S. Wichmann, \emph{Popoluca de Texistepec} (El Colegio de M\'exico, 2007; grammar, texts, lexicon, with CD). No direct URL to the book was found; a related overview: \url{https://arqueologiamexicana.mx/sites/default/files/banco_imagenes/popoluca-de-texistepec.pdf}
\item Texistepec language (Wikipedia): \url{https://en.wikipedia.org/wiki/Texistepec_language}
\item Glottolog entry: \url{https://glottolog.org/resource/languoid/id/texi1237}
\item Omniglot page on Texistepec: \url{https://www.omniglot.com/writing/texistepec.htm}
\item UNESCO Atlas entry: \url{https://ar.wal.unesco.org/index.php/en/countries/mexico/languages/texistepec-popoluca}
\end{itemize}
\section{Religious and Community Resources}
None found.
\section{Digital and Audio Resources}
\begin{itemize}
\item AILLA: Texistepec Zoque Collection of S{\o}ren Wichmann (audio of traditional stories): \url{https://islandora-ailla.lib.utexas.edu/islandora/object/ailla%3A124443}
\end{itemize}
\section{Curriculum Development Context}
Materials are research-oriented. A curriculum would build on Wichmann's grammar, texts and recordings together with the few remaining speakers.
\end{document}